% Storyline 1 (final), paragraph 2: the punishment dilemma and the two steps at which remedies can act.
The difficulty is sharpest for the norms that sustain cooperation under public information. Under these norms, a discriminator who refuses to help an individual in bad standing keeps its own good standing, because the refusal counts as justified punishment\cite{panchanathan2003,ohtsuki2004,okada2020dilemmas}. This exemption works only if the observer shares the donor's view of the recipient, which under private assessment is often not the case. An observer who still regards the recipient as good sees an unjustified refusal and condemns the discriminator. Disagreement about one individual thereby becomes disagreement about those who act towards it, and it spreads through the population\cite{hilbe2018,uchida2013,radzvilavicius2019}. Discriminators thus face a punishment dilemma. Withholding help from defectors protects cooperation against exploitation, yet each act of punishment risks their own reputation in the eyes of observers who disagree\cite{okada2020dilemmas,michelmata2024,pires2025}. The harm done by such disputed judgements depends on how often they occur and on how strongly each one changes a reputation. Remedies can therefore act at either step of reputation formation, on how actions are judged or on how judgements are turned into reputations, provided that genuine defectors still lose their standing.
